
\documentclass[10pt,aspectratio=169]{beamer}
\usepackage[utf8]{inputenc}
\usepackage[T1]{fontenc}
\usepackage[catalan]{babel}
\usepackage{amsmath,amssymb,booktabs,array,multicol}
\usepackage{tikz,pgfplots,xcolor,listings}
\usetikzlibrary{positioning,arrows.meta,calc,shapes.geometric}
\pgfplotsset{compat=1.18}
\usepgfplotslibrary{fillbetween}

\definecolor{UABGreen}{RGB}{0,103,71}
\definecolor{UABLight}{RGB}{224,241,236}
\definecolor{DeepBlue}{RGB}{20,65,110}
\definecolor{AccentOrange}{RGB}{230,126,34}
\definecolor{SoftGray}{RGB}{245,247,250}
\definecolor{CodeBack}{RGB}{248,248,248}
\definecolor{CodeFrame}{RGB}{210,215,220}

\usetheme{Madrid}
\usecolortheme[named=UABGreen]{structure}
\setbeamercolor{title}{fg=white,bg=UABGreen}
\setbeamercolor{frametitle}{fg=white,bg=UABGreen}
\setbeamercolor{block title}{fg=white,bg=DeepBlue}
\setbeamercolor{block body}{fg=black,bg=SoftGray}
\setbeamercolor{alerted text}{fg=AccentOrange}
\setbeamertemplate{navigation symbols}{}
\setbeamertemplate{itemize item}{\color{UABGreen}$\blacktriangleright$}
\setbeamertemplate{itemize subitem}{\color{AccentOrange}$\bullet$}
\setbeamertemplate{enumerate item}{\color{UABGreen}\insertenumlabel.}
\setbeamertemplate{footline}{%
  \leavevmode%
  \hbox{%
  \begin{beamercolorbox}[wd=.72\paperwidth,ht=2.4ex,dp=1ex,leftskip=1.4em]{author in head/foot}%
    Estadística -- Grau en Enginyeria Informàtica
  \end{beamercolorbox}%
  \begin{beamercolorbox}[wd=.28\paperwidth,ht=2.4ex,dp=1ex,center]{date in head/foot}%
    \insertframenumber{} / \inserttotalframenumber
  \end{beamercolorbox}}%
  \vskip0pt%
}

\lstdefinelanguage{R}{%
  morekeywords={if,else,repeat,while,function,for,in,next,break,TRUE,FALSE,NULL,NA,NaN,Inf,
    library,mean,median,sd,var,summary,table,prop.table,xtabs,addmargins,round,
    ggplot,aes,geom_point,geom_line,geom_col,geom_histogram,geom_density,geom_smooth,
    geom_tile,facet_wrap,labs,theme_minimal,summarise,group_by,mutate,read_csv,filter,
    arrange,set.seed,rbinom,rgeom,rpois,runif,rexp,rnorm,dnorm,pnorm,qnorm,dt,pt,qt,
    dchisq,pchisq,qchisq,df,pf,qf,dbinom,pbinom,dpois,ppois,replicate,sample,rep,seq,
    tibble,data.frame,confint,t.test,var.test,chisq.test,binom.test,prop.test,lm,
    broom,glance,augment,tidy,future_map_dbl,plan,multisession,map_dbl,purrr,future,furrr},
  sensitive=true,
  morecomment=[l]\#,
  morestring=[b]",
  morestring=[b]'
}
\lstset{language=R,basicstyle=\ttfamily\scriptsize,keywordstyle=\color{DeepBlue}\bfseries,
  commentstyle=\color{gray},stringstyle=\color{AccentOrange},backgroundcolor=\color{CodeBack},
  frame=single,rulecolor=\color{CodeFrame},breaklines=true,columns=fullflexible,keepspaces=true,
  showstringspaces=false,
  literate={à}{{\`a}}1 {è}{{\`e}}1 {é}{{\'e}}1 {í}{{\'i}}1 {ï}{{\"i}}1 {ò}{{\`o}}1 {ó}{{\'o}}1 {ú}{{\'u}}1 {ü}{{\"u}}1 {ç}{{\c{c}}}1 {À}{{\`A}}1 {È}{{\`E}}1 {É}{{\'E}}1 {Í}{{\'I}}1 {Ò}{{\`O}}1 {Ó}{{\'O}}1 {Ú}{{\'U}}1}

\newcommand{\R}{\textsf{R}}
\newcommand{\RStudio}{\textsf{RStudio/Posit}}
\newcommand{\GitHub}{\textsf{GitHub}}
\newcommand{\important}[1]{\textcolor{UABGreen}{\textbf{#1}}}
\newcommand{\exemple}[1]{\textcolor{AccentOrange}{\textbf{#1}}}
\newcommand{\Prob}{\mathbb{P}}
\newcommand{\E}{\mathbb{E}}
\newcommand{\Var}{\operatorname{Var}}
\newcommand{\IC}{\operatorname{IC}}

\title[Inferència I]{Inferència I}
\author{Estadística -- Grau en Enginyeria Informàtica}
\date{Curs 2026--27}

\begin{document}
\begin{frame}[plain]
  \titlepage
  \vspace{-0.45cm}
  \begin{center}
  \small\textcolor{DeepBlue}{Estadística computacional amb \R, simulació i exemples d'enginyeria informàtica}
  \end{center}
\end{frame}


\begin{frame}[fragile]{De la mostra a la decisió}
\begin{columns}[T]
\column{0.50\textwidth}
\begin{itemize}
  \item En descriptiva resumim dades observades.
  \item En inferència volem dir alguna cosa sobre una població o procés.
  \item La mostra és només una finestra parcial i aleatòria.
\end{itemize}
\column{0.44\textwidth}
\begin{block}{Exemples informàtics}
\begin{itemize}
  \item Temps mitjà de resposta d'un servidor.
  \item Proporció de builds que fallen.
  \item Variabilitat de potència en dispositius.
  \item Diferència de rendiment entre dues versions.
\end{itemize}
\end{block}
\end{columns}
\end{frame}

\begin{frame}[fragile]{Objectius del bloc d'inferència}
\begin{enumerate}
  \item Distingir població, mostra, paràmetre i estadístic.
  \item Construir intervals de confiança per quantificar incertesa.
  \item Interpretar correctament el nivell de confiança.
  \item Connectar fórmules amb càlculs en \R.
  \item Preparar el terreny per als tests d'hipòtesi.
\end{enumerate}
\end{frame}

\begin{frame}[fragile]{Població, mostra i estimadors}
\begin{columns}[T]
\column{0.50\textwidth}
\begin{block}{Població}
Procés o conjunt complet que ens interessa.
\[
\mu,\quad \sigma^2,\quad p
\]
\end{block}
\begin{block}{Mostra}
Subconjunt observat o repeticions mesurades.
\[
\bar X,\quad S^2,\quad \hat p
\]
\end{block}
\column{0.44\textwidth}
\begin{alertblock}{Exemple}
Volem estimar la latència mitjana real d'una API. Mesurem 200 peticions i obtenim $\bar X$ i $S$.
\end{alertblock}
\end{columns}
\end{frame}

\begin{frame}[fragile]{Interval de confiança: idea}
\begin{block}{Forma general}
\[
\text{estimador} \pm \text{marge d'error}.
\]
\end{block}
\begin{itemize}
  \item El marge d'error depèn de la variabilitat, la mida mostral i el nivell de confiança.
  \item Més dades $\Rightarrow$ interval més estret.
  \item Més confiança $\Rightarrow$ interval més ample.
\end{itemize}
\begin{alertblock}{Interpretació correcta}
Un 95\% de confiança vol dir que, si repetíssim el procediment moltes vegades, aproximadament el 95\% dels intervals construïts contindrien el paràmetre real.
\end{alertblock}
\end{frame}

\begin{frame}[fragile]{IC per a la mitjana amb $\sigma$ coneguda}
\begin{block}{Condicions}
$X$ normal o bé mida mostral prou gran.
\end{block}
\[
Z=\frac{\bar X-\mu}{\sigma/\sqrt n}\sim N(0,1)
\]
\[
\IC_{1-\alpha}(\mu)=
\bar X\pm z_{1-\alpha/2}\frac{\sigma}{\sqrt n}.
\]
\begin{lstlisting}
xbar <- 27.2; sigma <- 10.65; n <- 50
z <- qnorm(0.975)
e <- z * sigma / sqrt(n)
c(xbar - e, xbar + e)
\end{lstlisting}
\end{frame}

\begin{frame}[fragile]{IC per a la mitjana amb $\sigma$ desconeguda}
\begin{block}{Substituïm $\sigma$ per $S$}
\[
T=\frac{\bar X-\mu}{S/\sqrt n}\sim t_{n-1}
\]
\[
\IC_{1-\alpha}(\mu)=
\bar X\pm t_{n-1,1-\alpha/2}\frac{S}{\sqrt n}.
\]
\end{block}
\begin{lstlisting}
x <- c(115, 122, 119, 131, 125, 118, 121, 130)
t.test(x, conf.level = 0.95)$conf.int
\end{lstlisting}
\end{frame}

\begin{frame}[fragile]{Exemple computacional: latència mitjana}
\begin{columns}[T]
\column{0.50\textwidth}
Hem mesurat 80 peticions d'una API:
\begin{itemize}
  \item mitjana mostral: 118 ms,
  \item desviació típica mostral: 22 ms.
\end{itemize}
Volem un IC del 95\% per a la latència mitjana real.
\column{0.44\textwidth}
\begin{lstlisting}
xbar <- 118; s <- 22; n <- 80
crit <- qt(0.975, df = n - 1)
e <- crit * s / sqrt(n)
c(xbar - e, xbar + e)
\end{lstlisting}
\begin{block}{Lectura}
El resultat no és ``la latència sempre estarà dins l'interval'', sinó incertesa sobre la mitjana poblacional.
\end{block}
\end{columns}
\end{frame}

\begin{frame}[fragile]{Disseny mostral: quantes observacions necessitem?}
\begin{block}{Amb $\sigma$ aproximadament coneguda}
Per obtenir marge d'error màxim $e$:
\[
n\ge \left(\frac{z_{1-\alpha/2}\sigma}{e}\right)^2.
\]
\end{block}
\begin{lstlisting}
sigma <- 20     # ms
error <- 5      # marge màxim
z <- qnorm(0.975)
ceiling((z * sigma / error)^2)
\end{lstlisting}
\begin{itemize}
  \item En benchmarking, això ajuda a planificar quantes repeticions cal fer.
\end{itemize}
\end{frame}

\begin{frame}[fragile]{Flux de treball reproduïble}
\begin{enumerate}
  \item Definir la pregunta: què volem estimar?
  \item Recollir o simular dades de manera transparent.
  \item Calcular estimadors i interval.
  \item Visualitzar la mostra i discutir condicions.
  \item Escriure conclusions en un informe Quarto.
  \item Guardar codi i dades a \GitHub.
\end{enumerate}
\end{frame}

\end{document}
